
\documentclass[journal,transmag]{IEEEtran}
%
% If IEEEtran.cls has not been installed into the LaTeX system files,
% manually specify the path to it like:
% \documentclass[journal]{../sty/IEEEtran}



\usepackage{qtree} % just for figure example
\usepackage{listings} % for using listings...
\usepackage{xcolor}
\usepackage{amsfonts}
\usepackage{graphicx}
\usepackage{mathrsfs}
\usepackage{amsmath}%equation mutilrow align
\usepackage{amsthm}
\usepackage{amscd}
\usepackage{tabularx}
\usepackage{url}
\usepackage{multirow}
\usepackage{subfigure}
\usepackage{float}
\usepackage[caption=false]{subfig}
\usepackage{array}
\usepackage{algorithm}
\usepackage{amsmath}
\usepackage{bm}
\usepackage{comment}
\usepackage{makecell}
\newcommand*{\defeq}{\stackrel{\text{def}}{=}}
\usepackage[justification=centering]{caption}


%\usepackage[noend]{algpseudocode}

\usepackage{algpseudocode}

\algnewcommand\algorithmicforeach{\textbf{for each}}
\algdef{S}[FOR]{ForEach}[1]{\algorithmicforeach\ #1\ \algorithmicdo}

\newtheorem{definition}{Definition}
\newtheorem{theorem}{Theorem}%[section]
\newtheorem{lemma}{Lemma}
\newtheorem{proposition}{Proposition}
\newtheorem{corollary}{Corollary}
\newtheorem{assumption}{Assumption}

\setcounter{page}{1}






% *** GRAPHICS RELATED PACKAGES ***
%
\ifCLASSINFOpdf
  % \usepackage[pdftex]{graphicx}
  % declare the path(s) where your graphic files are
  % \graphicspath{{../pdf/}{../jpeg/}}
  % and their extensions so you won't have to specify these with
  % every instance of \includegraphics
  % \DeclareGraphicsExtensions{.pdf,.jpeg,.png}
\else
  % or other class option (dvipsone, dvipdf, if not using dvips). graphicx
  % will default to the driver specified in the system graphics.cfg if no
  % driver is specified.
  % \usepackage[dvips]{graphicx}
  % declare the path(s) where your graphic files are
  % \graphicspath{{../eps/}}
  % and their extensions so you won't have to specify these with
  % every instance of \includegraphics
  % \DeclareGraphicsExtensions{.eps}
\fi


% correct bad hyphenation here
\hyphenation{op-tical net-works semi-conduc-tor} 




\begin{document}
%
% paper title
% can use linebreaks \\ within to get better formatting as desired
% Do not put math or special symbols in the title.
\title{Appendix}

% author names and affiliations
% transmag papers use the long conference author name format.



\begin{comment}
\author{\IEEEauthorblockN{Bo Ma\IEEEauthorrefmark{1}
%,
%Jinsong Wu\IEEEauthorrefmark{3}}
}
\IEEEauthorblockA{\IEEEauthorrefmark{1}School of Mathematics and Computer Engineering,
Auckland University of Technology,Auckalnd  1024 NZ}

 % stops an unwanted space
\thanks{Manuscript received July 1, 2020; revised July 27, 2020.
Corresponding author:Bo Ma@aut.ac.nz}}
\end{comment}


% The paper headers
%\markboth{IEEE Transaction on Pattern Analysis and Machine Intelligence ,~Vol.~11, No.~4, March~2021}%
%{Shell \MakeLowercase{\textit{et al.}}: Bo Ma of IEEEtran.cls for Journals}




% make the title area
\maketitle


% To allow for easy dual compilation without having to reenter the
% abstract/keywords data, the \IEEEtitleabstractindextext text will
% not be used in maketitle, but will appear (i.e., to be "transported")
% here as \IEEEdisplaynontitleabstractindextext when the compsoc
% or transmag modes are not selected <OR> if conference mode is selected
% - because all conference papers position the abstract like regular
% papers do.
\IEEEdisplaynontitleabstractindextext
% \IEEEdisplaynontitleabstractindextext has no effect when using
% compsoc or transmag under a non-conference mode.







% For peer review papers, you can put extra information on the cover
% page as needed:
% \ifCLASSOPTIONpeerreview
% \begin{center} \bfseries EDICS Category: 3-BBND \end{center}
% \fi
%
% For peerreview papers, this IEEEtran command inserts a page break and
% creates the second title. It will be ignored for other modes.
%\IEEEpeerreviewmaketitle


\section*{Appendix}
\subsection{Related Works}
The main objective is to detect the target objects while ensuring privacy.  
Several deep learning methods have achieved very high accuracy and efficiency in detecting objects in images. 
Among them, Fast R-CNN~\cite{ren2016faster} and Mask R-CNN~\cite{he2017mask} extract the features of the entire image at the same time.
These features are then passed to the spatial pyramid pool and the interest domain pool layer to generate regional features. 
Region of Interest (ROI) pooling extends the ROI align layer to reduce the error of spatial quantization~\cite{puzicha2000spatial}. 
These are followed by Faster R-CNN~\cite{ren2016faster} which is more computationally efficient than the previous two. 
The Feature Pyramid Network (FPN)~\cite{lin2017feature} is later proposed which uses the Faster R-CNN as its backbone network to replace the independent and time-consuming area select method.

In terms of speed and accuracy, the You Only Look Once (YOLO) model~\cite{farhadi2018yolov3,redmon2016you} is one of the latest deep neural network object detection methods. 
YOLO is able to perform object detection for videos in embedded devices in real-time. Furthermore, it uses the evolved deep intelligence framework to evolve the YOLO network architecture and produce an optimized architecture.
Like YOLO, SSD~\cite{liu2016ssd} is also a one-stage and end-to-end object detection method. 
It directly classifies predefined anchors and uses convolutional neural networks for detection. 
Yolo-V5 has incorporated the multilayer prediction function found in SSD and introduced additional context information with deconvolution operators to improve accuracy.

Inspired by the success of using multi-headed attention in Transformer architectures for natural language processing (MLP), many researchers have proposed variants of the Transformer for image classification and object detection.
The Detection Transformer (DETR)~\cite{carion2020end} performs end-to-end object detection. 
It trains data with a bipartite matching loss and is easily reproduced since it does not have any customized layers.
The Vision Transformer (ViT)~\cite{dosovitskiy2020image} interprets an image as a sequence of patches.
It applies the encoder of the standard transformer with an additional learnable classification token~\cite{vaswani2017attention} to process the input. 
To compensate for the weakness of ViT in performing pixel-level dense detection, the Pyramid Vision Transformer (PVT)~\cite{wang2021pyramid} was proposed. It can effectively detect dense images by replacing the backbone of CNN.
Fine-grained image patches were used to provide high-resolution representation. 
In contrast, in order to operate under limited resources, both the progressive shrinking pyramid and spatial-reduction attention were incorporated to acquire multi-scale feature maps. 
With convolutional backbones, both feature resolution and granularity are compromised, which results in difficulties in recovering the decoders. Dense prediction transformer (DPT)~\cite{ranftl2021vision} utilizes the ViT and convolutional network as encoder and decoder respectively. Noticeably, per-stage ViT encoders retain spatial resolution within the original embedding. 

The use of attention modules is also inspired by the Transformer architecture.
There are generally two ways to apply the attention mechanism.
This first one is per-channel attention and the second is a combination of channel and spatial attention. Squeeze-and-Excite (SE)~\cite{hu2018squeeze} shows excellent potential in efficiency through dimensionality reduction. However, capturing dependencies across all channels tends to increase computational complexity significantly. 
Efficient Channel Attention (ECA)~\cite{qilong2020eca} acquires interaction across channels through a lightweight method. Nevertheless, both of these methods disregard the spatial dimension. 
The Bottleneck Attention Module (BAM)~\cite{park2018bam} and Convolutional Block Attention Module (CBAM)~\cite{woo2018cbam} adopt channel and spatial attention to refine convolutional features. 
While the transformer can detect the object and classify the images according to its feature, protecting privacy according to a privacy budget is still a problem. The reason is that it is not easy to fit differential privacy sub-sampling with the Transformer framework.
% Can use something like this to put references on a page
% by themselves when using endfloat and the captionsoff option.
\ifCLASSOPTIONcaptionsoff
  \newpage
\fi


% trigger a \newpage just before the given reference
% number - used to balance the columns on the last page
% adjust value as needed - may need to be readjusted if
% the document is modified later
%\IEEEtriggeratref{8}
% The "triggered" command can be changed if desired:
%\IEEEtriggercmd{\enlargethispage{-5in}}

% references section

% can use a bibliography generated by BibTeX as a .bbl file
% BibTeX documentation can be easily obtained at:
% http://www.ctan.org/tex-archive/biblio/bibtex/contrib/doc/
% The IEEEtran BibTeX style support page is at:
% http://www.michaelshell.org/tex/ieeetran/bibtex/
%\bibliographystyle{IEEEtran}
% argument is your BibTeX string definitions and bibliography database(s)
%\bibliography{IEEEabrv,../bib/paper}
%
% <OR> manually copy in the resultant .bbl file
% set second argument of \begin to the number of references
% (used to reserve space for the reference number labels box)
\bibliographystyle{IEEEtran}
\bibliography{IEEEabrv,ref}

% biography section
%
% If you have an EPS/PDF photo (graphicx package needed) extra braces are
% needed around the contents of the optional argument to biography to prevent
% the LaTeX parser from getting confused when it sees the complicated
% \includegraphics command within an optional argument. (You could create
% your own custom macro containing the \includegraphics command to make things
% simpler here.)
%\begin{IEEEbiography}[{\includegraphics[width=1in,height=1.25in,clip,keepaspectratio]{mshell}}]{Michael Shell}
% or if you just want to reserve a space for a photo:
\begin{comment}
\begin{IEEEbiography}
[{\includegraphics[width=1in,height=1.25in,clip,keepaspectratio]{images/boma.jpg}}]{Bo Ma}
Ma is working toward the PhD degree in the Department of Information Technology and Software Engineering, School of Engineering, Computer and Mathematical Sciences at the Auckland University of Technology, Auckland, New Zealand. His research interest is Cybersecurity for various networks and systems in general, specially focus on privacy-preserving federated machine learning methods, privacy issues of Big Data analytics and Internet of Things, and information theory of deep learning. He obtained his MSc in Computer Software and Theory, and
Bachelor in Electronic Commerce from Sichuan Normal University, China.
\end{IEEEbiography}



% if you will not have a photo at all:
\begin{IEEEbiographynophoto}{John Doe}
Biography text here.
\end{IEEEbiographynophoto}

% insert where needed to balance the two columns on the last page with
% biographies
%\newpage

\begin{IEEEbiographynophoto}{Jane Doe}
Biography text here.
\end{IEEEbiographynophoto}
\end{comment}
% You can push biographies down or up by placing
% a \vfill before or after them. The appropriate
% use of \vfill depends on what kind of text is
% on the last page and whether or not the columns
% are being equalized.

%\vfill

% Can be used to pull up biographies so that the bottom of the last one
% is flush with the other column.
%\enlargethispage{-5in}



% that's all folks
\end{document}

